\chapter{Glossary of Industry Terms}
\label{app:glossary}
\addcontentsline{toc}{chapter}{Appendix E: Glossary of Industry Terms}

\noindent
This glossary defines the industry-specific terms used in the
thesis.  Definitions follow the operational usage of the
Tamil-cinema and Indian-OTT industries circa 2024--2025; where the
academic literature uses a different operationalisation, the
academic source is cited.

\begin{description}[style=nextline,leftmargin=2.4em]
\item[A-list star.]  An actor whose presence in a Tamil-language
        feature film commands an opening-weekend audience pull
        materially above the median, but who does not belong to the
        small ``top-tier'' set of four to five actors whose name
        alone can carry a film.  In specialist usage in 2024, the
        Tamil A-list comprises approximately fifteen male and seven
        female actors.

\item[Acquisition price.]  The lump-sum amount paid by an OTT
        platform to a producer to acquire the streaming rights of
        a Tamil-language feature film.  Reported either in absolute
        rupees or as an index against a 2018 baseline
        (Section~\ref{sec:price_index}).

\item[Audio launch.]  The public release event for a Tamil-language
        film's soundtrack, typically held two to four months before
        theatrical release, attended by 8{,}000--40{,}000 invitees.

\item[AVOD.]  Advertising-supported Video On Demand.  A streaming
        service whose monetisation model is advertising rather than
        subscription, typically free at the point of consumption.

\item[B-list star.]  An actor below the A-list whose presence
        contributes meaningfully to opening-day pull but does not by
        itself sustain a film's commercial proposition.

\item[Box-office gross.]  The total ticketed revenue collected
        across all theatrical screenings of a film.

\item[Direct-to-OTT.]  A film whose first paid public exhibition
        occurs on a streaming platform rather than in a theatre.

\item[Distributor.]  A firm or individual that acquires the
        territorial right to exhibit a film in specified circuits
        and shares revenue with the producer on agreed terms.

\item[Exhibition geography.]  The spatial distribution of single-screen
        and multiplex cinemas in a given market, typically reported
        at the state, district or town level.

\item[FDFS.]  First Day First Show -- the first public screening of
        a star-led Tamil film, traditionally held at 6 am on a
        Friday and accompanied by fan-club mobilisation, milk
        ablutions of star portraits and ceremonial garlanding.

\item[Greenlight.]  The formal commitment by a producer to fund and
        commence production of a film.

\item[Hybrid window.]  The 3--4 week theatre-to-OTT window that has
        emerged as the modal Tamil-cinema arrangement post-2022
        (Section~\ref{sec:windowing_findings}).

\item[Multiplex.]  A cinema with three or more screens under unified
        operation, typically equipped with digital projection,
        air-conditioning, premium seating and recognisable brand
        operation (PVR-Inox, SPI / Sathyam, AGS, Cinepolis, Rohini
        and others).

\item[OTT.]  Over-the-Top -- an audiovisual service delivered
        directly to internet-connected devices, bypassing
        traditional broadcast or theatrical distribution
        \citep{lobato2019}.

\item[Pan-India release.]  A theatrical release strategy in which a
        film is released simultaneously in two or more Indian
        languages -- typically Tamil, Telugu, Hindi, Malayalam and
        Kannada -- with synchronised marketing and exhibition.

\item[Pre-sale economics.]  The financing model in which a film is
        partially or fully funded by selling its territorial,
        satellite and audio rights ahead of theatrical release.

\item[Producer.]  The firm or individual that holds the financial
        and creative authority over a film's production.

\item[Production-cultures research.]  An ethnographic-and-textual
        approach to studying media industries that focuses on how
        industry workers themselves theorise and narrate change
        \citep{caldwell2008}.

\item[Rasikar mandram.]  Tamil for ``fan club''; the
        hierarchically-organised infrastructure of fan groups
        attached to a major Tamil star.

\item[Re-windowing.]  The re-engineering of historical theatrical /
        home-video / pay-TV / free-TV exploitation sequences in the
        streaming era \citep{cunningham2019}.

\item[Satellite rights.]  The right to broadcast a film on
        satellite television, typically sold by the producer to a
        broadcaster (Sun TV, KTV, Vijay TV, Zee Tamil) for a
        lump-sum fee within an agreed-on post-theatrical window.

\item[Single-screen.]  A cinema operating one auditorium, typically
        of older vintage, often (in Tamil Nadu) family-owned and
        located in Tier-2 / Tier-3 towns.

\item[Subscription Video On Demand (SVOD).]  A streaming service
        whose primary monetisation is recurring user subscription
        (Netflix, Amazon Prime Video, Disney+ Hotstar, Sun NXT,
        Aha Tamil, ZEE5, JioCinema premium tier, etc.).

\item[Star-fee adjustment.]  The percentage change in fees commanded
        by an actor between a baseline year (2019) and a comparison
        year (2024); reported by tier in
        Section~\ref{sec:star_fees_results}.

\item[Top-tier star.]  An actor whose presence in a Tamil-language
        feature film independently sustains the film's commercial
        proposition.  In contemporary Tamil cinema, top-tier is
        commonly identified with Vijay, Ajith, Kamal Haasan,
        Rajinikanth and -- depending on the analyst -- a small
        number of additional names.

\item[Theatrical window.]  The number of days between a film's
        first paid theatrical screening and its first paid digital /
        OTT availability \citep{vogel2020}.

\item[TVOD.]  Transactional Video On Demand -- a streaming service
        in which the user pays per-title rather than via
        subscription, typical of Apple iTunes-style rentals.

\item[Word-of-mouth (WOM).]  The informal post-release transmission
        of audience verdict via interpersonal communication, today
        increasingly mediated by WhatsApp, Twitter / X, Instagram
        Reels and YouTube reaction videos.
\end{description}
